\documentclass[conference]{IEEEtran}
\IEEEoverridecommandlockouts
% The preceding line is only needed to identify funding in the first footnote. If that is unneeded, please comment it out.
\usepackage{cite}
\usepackage{amsmath,amssymb,amsfonts}
\usepackage{algorithmic}
\usepackage{graphicx}
\usepackage{textcomp}
\usepackage{xcolor}
\def\BibTeX{{\rm B\kern-.05em{\sc i\kern-.025em b}\kern-.08em
    T\kern-.1667em\lower.7ex\hbox{E}\kern-.125emX}}
\begin{document}

\title{An Explainable Architecture for Career Recommendation and Skill-Gap-Aware Learning Path Generation}

\author{\IEEEauthorblockN{First Author}
\IEEEauthorblockA{\textit{Department of Computer Science and Engineering} \\
\textit{GHRCE}\\
Nagpur, India \\
first.author@example.com}
\and
\IEEEauthorblockN{Second Author}
\IEEEauthorblockA{\textit{Centre of Excellence -- Data Science} \\
\textit{GHRCE}\\
Nagpur, India \\
second.author@example.com}
}

\maketitle

\begin{abstract}
As the rapid evolution of technology complicates engineering career trajectories, students require reliable mechanisms to align their academic skills with industry roles. Current recommendations systems increasingly rely on deep-learning models or Large Language Models (LLMs), which demand significant computational resources, introduce hallucination risks, and lack mathematical explainability for critical career mapping. This paper proposes PATHFORGE, a hybrid framework that isolates deterministic skill analysis from LLM-driven conversational mentorship. By employing a lightweight, heuristic-based scoring engine (BYSER), the system maps user competencies to career requirements securely and sequentially generates static, dependency-graph-based learning roadmaps. A multi-agent LLM layer subsequently utilizes these mathematically verified outputs as contextual grounding to orchestrate conversational mentors. Simulated benchmarking of our deterministic pipeline (N= 200 student profiles) reveals a top-1 domain alignment accuracy of 100\% and an average computation latency of 0.037 ms per match, validating its efficacy as an explainable, scalable decision-support tool. 
\end{abstract}

\begin{IEEEkeywords}
Career Recommendation, Skill Gap Analysis, Explainable AI, Personalized Learning, LLM Orchestration
\end{IEEEkeywords}

\section{Introduction}
The diversification of software engineering and artificial intelligence disciplines has created a significant guidance gap for engineering undergraduates attempting to select specialization paths. The challenge is twofold: accurately identifying overlapping skills between a student's profile and dynamic market requirements, and formulating a sequential learning roadmap to bridge isolated skill gaps \cite{prior1, prior2}.

Recent advancements in recommender systems increasingly deploy deep neural networks and transformer-based semantic retrieval to map user resumes against employment listings \cite{semantic1}. Concurrently, Large Language Models (LLMs) are engaged to dynamically generate personalized study schedules \cite{llm1}. While powerful, these stochastic methodologies suffer from critical application limits within academic environments. Deep learning vectors are highly uninterpretable, making it difficult to explicitly justify to a student \textit{why} a career was chosen \cite{xai1}. Furthermore, delegating sequential syllabus creation to probabilistic LLMs frequently results in hallucinatory prerequisite structures (e.g., advising studying React.js before JavaScript), undermining educational integrity.

To address these limitations, we present a novel hybrid system architecture that rigorously separates deterministic skill evaluation from generative conversational oversight. The proposed system features the BYSER matching engine, which operates via an explainable mathematical heuristic across numerical attributes (skills, academics, coding, projects) rather than hidden embedding weights. Roadmaps are not hallucinated; they are dynamically traversed from a static topological skill dependency graph. LLMs (specifically Google Gemini 1.5) are then restricted solely to conversational roles, using the deterministic outputs of the system as their systemic constraints. 

This paper details the exact implementation of this architecture, provides extensive methodology surrounding its rule-based generation systems, and presents benchmark computations validating its applicability. 

\section{Related Work}

\subsection{Career Prediction and Semantic Matching}
Early computational career prediction utilized traditional classifiers such as Random Forests and Support Vector Machines, reliant on large sets of historical user data \cite{prior1}. As transformers gained prominence, research transitioned to BERT-based semantic extraction, evaluating cosine similarities between a candidate's resume embeddings and standard occupational matrices such as O*NET and ESCO \cite{semantic1, esco1}. Although semantic matching easily captures synonymous skill variants, the transformation of explicit student competencies into opaque numerical vectors sacrifices interpretability. 

\subsection{Personalized Roadmap Generation and LLMs}
In parallel, the education domain has explored personalized roadmap generation using Item Response Theory and collaborative filtering. Modern implementations leverage LLMs for conversational agents (RAG architecture) that provide dynamic mentorship \cite{llm1}. However, studies evaluating LLM autonomy highlight significant risks associated with automation bias and curriculum hallucination when LLMs orchestrate multi-month dependencies independently \cite{xai1}. 

\subsection{Positioning of the Proposed System}
Differing from prior work that combines career matching and learning path generation directly inside implicit predictive models, PATHFORGE positions the LLM strictly as a human-computer interaction layer. It substitutes implicit neural-recommendations with an explicit algebraic framework, prioritizing transparency, reliability, and extreme computational efficiency over probabilistic generalization.

\section{Proposed System Architecture and Methodology}

The implemented architecture is a multi-tier monolithic application communicating over RESTful API interfaces utilizing React, Express.js, and PostgreSQL.

\subsection{User and Career Representation}
Instead of generating mathematical embeddings natively in production, both Students ($U$) and Careers ($C$) are represented as standardized JSON arrays tied strictly to their respective explicit taxonomies. For an arbitrary user profile $U$, let $S_u$ denote a set of possessed skills, where each skill $s \in S_u$ possesses an integer proficiency grading, $P(s) \in \{1, 2, 3, 4, 5\}$. A candidate career $C$ enforces required skills $S_c$, each coupled to an assigned importance weight $W_c(s) \in [1, 10]$.

\subsection{The BYSER Career Recommendation Engine}
Unlike typical collaborative or content-based machine learning approaches, career recommendation is orchestrated securely through a static formula designated BYSER. The engine calculates an explicit matching score across 7 distinct dimensions: Skills (35\%), Interests (20\%), Projects (15\%), Academics (10\%), Certifications (10\%), Coding Experience (5\%), and Career Goals (5\%).

The foundational Skills Metric ($SM$) is formalized as the ratio of possessed skill weight to required proficiency:
\begin{equation}
SM_{u,c} = \frac{\sum_{s \in (S_u \cap S_c)} \left( W_c(s) \times \frac{P(s)}{5} \right)}{\sum_{s \in S_c} W_c(s)} \times 100
\end{equation}

Categorical attributes, such as project experience match, are evaluated via localized string evaluations traversing the target career categorization, acting as discrete bonus thresholds. The exact numerical bounds dictate the Final Recommended Score $Total(U, C)$. Careers are ranked in descending order, ensuring absolute numerical explainability for every parameter contributing to the student's result.

\subsection{Skill-Gap Index (SGI)}
The Skill Gap Index corresponds directly towards the complement of the Skill Metric match.
\begin{equation}
SGI(U, C) = \max(0, 100 - SM_{u,c})
\end{equation}
Nodes missing inside $S_c$ which fail to exist inside the student repository $S_u$ are categorized dynamically as "High", "Medium," or "Low" priority based solely upon their predetermined relational weight $W_c(s)$.

\subsection{Personalized Roadmap Engine}
LLMs natively fail at graphing rigorous prerequisites over deeply linked networks. To rectify this, the roadmap logic is defined over a topological Directed Acyclic Graph (DAG), traversing standard curriculum prerequisites (e.g., HTML $\rightarrow$ JavaScript $\rightarrow$ TypeScript $\rightarrow$ Angular). For an identified $SGI_{u,c}$, the algorithmic iterator visits missing node definitions recursively to compile an ordered sequence array. 

This array is statically mapped across 3 learning phases (Fundamentals, Core Engineering, and Capstone Deployment) adjusted by a pace multiplier modifying chronological estimation, natively yielding a fully populated educational roadmap JSON structure populated without single AI API usage.

\subsection{Multi-Agent LLM Orchestration Integration}
Machine Intelligence represents the pinnacle UI. Using the Google GenAI SDK (Gemini 1.5 API), a suite of restricted intelligent agents (e.g., Interview Coach, Resume Builder, Research Mentor) operates atop an isolated orchestrator framework. Rather than tasking the LLM to process career decisions directly, the REST server automatically compiles the transparent BYSER ranking evaluation array, Skill Map dependencies, and User representations, explicitly formatting them as systemic text instructions injected securely inside the prompt context per Agent Request. This architecture entirely circumvents prompt context-rot and algorithm hallucination.

\section{Experimental Configuration and Results}

The core assertion dictates that heuristic logic significantly eclipses standard deep learning methodologies in throughput and resource efficiency whilst assuring deterministic matching reliability for student alignment. In the absence of available open-source ground truth student mappings specific to GHRCE internal metrics, we developed a synthetic computational benchmark mapping 200 mathematically disparate engineering profiles categorized identically across 3 strict domain orientations (Frontend, Backend, and AI Data Engineering) spanning random GPA distributions.

\subsection{Algorithmic Reliability}
Table \ref{tab:accuracy} details the deterministic alignment matching evaluation processed identically via the BYSER runtime loop. No profiles experienced heuristic collision mapping failures. The pipeline retained 100\% mathematical predictability ranking alignment.

\begin{table}[htbp]
\caption{System Routing Alignment Accuracy over Subsets}
\begin{center}
\begin{tabular}{|c|c|c|}
\hline
\textbf{Focus Distribution} & \textbf{Trials} & \textbf{Top-1 Routing Alignment} \\
\hline
Frontend Profiles & 67 & 100.00\% \\
Backend Profiles & 67 & 100.00\% \\
AI/ML Profiles & 66 & 100.00\% \\
\hline
\end{tabular}
\label{tab:accuracy}
\end{center}
\end{table}

\subsection{Inference Latency Profile}
To evaluate scalability feasibility within intensive concurrent request bursts inside academic server infrastructure, computational latency arrays were extracted natively over Node.js profiling.

\begin{table}[htbp]
\caption{Calculation Scalability Metrics (N=200)}
\begin{center}
\begin{tabular}{|c|c|}
\hline
\textbf{Metric Parameter} & \textbf{Recorded Data} \\
\hline
Total System Execution & 7.42 \text{ ms} \\
Average Latency Per Match & 0.037 \text{ ms} \\
\hline
\end{tabular}
\label{tab:latency}
\end{center}
\end{table}

Traditional Transformer-based classification deployments routinely process requests extending between 150ms to 400ms \cite{llm1}. The calculated BYSER sub-millisecond scaling threshold (0.037 ms/match) justifies avoiding massive compute deployments in favor of explainable relational database algorithms dynamically wrapping a frontend API interface. 

\section{Discussion and Limitations}

The empirical evaluation clarifies architecture robustness; however, strict exact-match methodologies possess innate weaknesses lacking semantic generalization techniques. 
\begin{itemize}
\item \textit{Exact String Coupling:} The heuristic algorithm operates via substring similarity matching (`.toLowerCase() ===`). A student mistakenly inputting "JS" instead of "JavaScript" would suffer zero points mapped in the DB relations.
\item\textit{External Data Provenance Limitations:} Placeholder static URL strings route dynamic learning materials; ensuring sustained relevance demands API implementation mapped to evolving repositories (e.g., Coursera schemas). 
\item \textit{Evaluation Limits:} Longitudinally assessing whether generated roadmaps demonstrably improve final semester job placements implies long-term human evaluation pending actual real-world institutional cohort tracking trials. 
\end{itemize}

\section{Conclusion}

This paper has demonstrated PATHFORGE, a highly transparent hybrid web framework dedicated toward explainable academic career alignment. The research challenges the industry dependence upon inherently opaque machine learning networks for personalized education. Implementing deterministic mathematical engines (BYSER) and graphed dependencies successfully segregates functional processing execution from AI conversational layers, significantly mitigating LLM risk. With computational matching latencies registering below 0.05 milliseconds and flawless deterministic scaling evaluation, the implemented blueprint validates lightweight computation mapping architectures optimized exactly for scalable academic mentorship systems.

\bibliographystyle{IEEEtran}
\bibliography{references}

\end{document}
